% Source: mathematics/phase_space_return_gas.md
% Source SHA256: 95acc1bc4b196a68852a62cc813201bf2494a9ad36083fc45ec58d7f4ead6385
\hypertarget{general-analytic-phase-space-kicks-and-the-regularity-of-return-averages}{%
\section{General analytic phase-space kicks and the regularity of return averages}\label{general-analytic-phase-space-kicks-and-the-regularity-of-return-averages}}
\par\smallskip\noindent\emph{Numbering convention: section and theorem references within this chapter use the source module\textquotesingle s local numbering; source Section 1 is subsection \thesection.1. Equation references carry the prefix Ret.}\par\smallskip

\textbf{Status.} Project derivation with a limited independent internal check
of the vector lattice, Weyl phases, analytic gas, and regularity proofs.
Exact finite arithmetic checks and floating matrix diagnostics are
separated in Section 10. External refereeing, priority verification,
and Lean formalization remain open. No microscopic BGS universality
is claimed.

Use the balanced quantization, arithmetic sequence and notation of
\href{https://github.com/ChurilovMaksim/BGS/blob/main/mathematics/periodic_point_trace_comparison.md}{the periodic-point comparison}:
\[
 A=\begin{pmatrix}3&4\\2&3\end{pmatrix},\quad T=A^{\mathsf T},
 \quad\lambda=3+2\sqrt2,\quad K=2m,\quad m\ge3\ {\rm odd},
 \quad N=b_m,\quad U_0^K=I .
\]
The scalar-site gas and its classical Penrose tree identity are proved in
\href{https://github.com/ChurilovMaksim/BGS/blob/main/mathematics/fixed_harmonic_return_gas.md}{the fixed-harmonic note}, Section 3.
The argument below retains the additional onsite and antipodal phases
that arise for a general two-dimensional symbol.

\hypertarget{analytic-return-theorem}{%
\subsection{Analytic return theorem}\label{analytic-return-theorem}}

Let \(f\) be a fixed real analytic mean-zero function on \(\mathbb T^2\).
Equivalently for our purposes, assume
\[
 F_\rho(f)=\sum_{v\in\mathbb Z^2}|\widehat f_v|e^{\rho\|v\|_1}
                                  <\infty
 \quad\hbox{for some }\rho>0,\qquad \widehat f_0=0.
\]
Let \(\sigma^2=\sigma_A^2(f)\) be the full orbit variance (Ret.P18).

\textbf{Theorem 1.} If \(\sigma^2>0\), there are \(s_0,C,c>0\), depending on
\(f\), such that for every admissible \(K\), every positive integer \(r\),
and every real \(\epsilon\) with \(s=r|\epsilon|\le s_0\),
\[
 \left|\frac1N\operatorname{Tr}U_{\epsilon,f}^{rK}\right|
                 \le e^{-\sigma^2Ks^2/4}+Ce^{-cK},          \tag{Ret.R1}
\]
\[
 \left|\frac1N\operatorname{Tr}U_{\epsilon,f}^{rK}
                          -e^{-\sigma^2Ks^2/2}\right|
                 \le\frac C{\sqrt K}+Ce^{-cK}.             \tag{Ret.R2}
\]
For even \(f\) the same estimates hold on the normalized odd block.
Constants are not uniform as \(f\) approaches the zero-variance subspace.
If \(\sigma^2=0\), the stronger coherent estimate is
\[
 \left|\frac1N\operatorname{Tr}U_{\epsilon,f}^{rK}-1\right|
                         \le \frac{C_fKs^2}{rN}.            \tag{Ret.R3}
\]

Thus for every \(\epsilon_K\to0\) and fixed positive integer \(r\),
the trace differs from \(e^{-\sigma^2r^2K\epsilon_K^2/2}\) by \(o(1)\).
When \(\sigma^2>0\), the global eigenphase law of \(U_{\epsilon_K,f}^K\)
is Haar measure on the circle if and only if \(K\epsilon_K^2\to\infty\).
For \(K\epsilon_K^2\to v<\infty\) it is the wrapped Gaussian with variance
\(v\sigma^2\). For \(\sigma^2=0\) the return operator tends to the identity
in operator norm. These are global return laws at vanishing classical
kick strength, not statements about local level spacings.

\hypertarget{analytic-cohomology-and-orbit-representatives}{%
\subsection{Analytic cohomology and orbit representatives}\label{analytic-cohomology-and-orbit-representatives}}

Choose the eigen-coordinate norm \(|v|_*\) of Section 10 of the companion
note. In each oriented nonzero \(A\)-orbit choose a representative of
minimal \(|\cdot|_*\), with paired representatives for negative orbits.
Put
\[
 b_{v_{\mathcal O}}=\sum_{v\in\mathcal O}\widehat f_v,\qquad
 b(z)=\sum_{\mathcal O}b_{v_{\mathcal O}}e_{v_{\mathcal O}}(z).
\]
Then there is a real analytic \(g\) such that
\[
 f=b+g-g\circ T,\qquad
 \int b^2=\sum_{\mathcal O}|b_{v_{\mathcal O}}|^2=\sigma^2.  \tag{Ret.R4}
\]
The nonzero frequencies of \(b\) meet each oriented \(A\)-orbit at most once.
Both \(b\) and \(g\) are even when \(f\) is even.

For completeness, analyticity is preserved with a positive strip width.
First use an exponential weight in the equivalent norm \(|\cdot|_*\).
Minimality gives \(|v_{\mathcal O}|_*\le|v|_*\) on its orbit, so the
weighted Fourier norm of \(b\) is at most that of \(f\).
Write \(v_j=A^jv_0\) on one orbit and define
\[
 \widehat g_{v_j}=
 \begin{cases}
   \sum_{\ell\le j}\widehat f_{v_\ell},&j<0,\\
   -\sum_{\ell>j}\widehat f_{v_\ell},&j\ge0.
 \end{cases}
\]
The difference at \(j=0\) is \(\widehat f_{v_0}-b_{v_0}\), as required.
On each half-orbit the norms grow geometrically from their minimum.
In particular \(|v_j|_*\le|v_\ell|_*\) between the representative and
\(v_\ell\), and \(|\ell|\le C_A\log(2+|v_\ell|_*)\).
For \(0<\rho'<\rho\), summing the transfer coefficients with weight
\(e^{\rho'|v|_*}\) is therefore bounded by
\[
 C_A\sum_v|\widehat f_v|
                  \log(2+|v|_*)e^{\rho'|v|_*}
 \le C_{A,\rho,\rho'}\sum_v|\widehat f_v|e^{\rho|v|_*}.
\]
This proves (Ret.R4) and analyticity after reverting to \(\|\cdot\|_1\).

Two-factor unitary Duhamel applied to
\(e^{i\epsilon\operatorname{Op}(g)}  e^{i\epsilon\operatorname{Op}(b)}  e^{-i\epsilon\operatorname{Op}(g\circ T)}\)
shows that \(U_{\epsilon,f}\) is unitarily conjugate to \(U_{\epsilon,b}\)
up to operator error \(C_f\epsilon^2/N\). Each commutator of two fixed
smooth symbols costs at most \(2\pi F_1(a)F_1(d)/N\).
Consequently the \(rK\)-trace error is at most
\(C_fKs^2/(rN)\), with the same result on an invariant odd block.
If \(\sigma^2=0\), all coefficients of \(b\) vanish and (Ret.R3) follows.
It remains to prove (Ret.R1)--(Ret.R2) for an analytic \(b\) with the stated
orbit-representative property and \(\int b^2>0\).

\hypertarget{the-entire-vector-resonance-lattice}{%
\subsection{The entire vector resonance lattice}\label{the-entire-vector-resonance-lattice}}

On cyclic vector sequences let
\[
 (\mathcal Dp)_j=p_j-Ap_{j+1},\qquad
 \mathcal C_K=\left\{c:\sum_jA^jc_j=0\pmod N\right\},
 \quad q(c)=|\{j:c_j\ne0\}|.
\]
The companion note proves
\(p=\mathcal D^{-1}c\in(1/32)\mathbb Z^{2K}\) and
\(\|p\|_1\le2\|c\|_1\) for \(c\in\mathcal C_K\).
If one \(p_j\) is an integer vector, integrality propagates around the
circle. Thus any noninteger class has no zero vector coordinate.

For a nonzero rational vector \(p=(x,y)\) of denominator dividing 32,
\[
 Q(p)=\omega(p,Ap)=2x^2-4y^2\ne0,\qquad |Q(p)|\ge1/1024.
\]
In coordinates \(p=\xi(\sqrt2,1)+\zeta(-\sqrt2,1)\),
\(Q=-16\xi\zeta\) and \(\|p\|_1\ge2\max(|\xi|,|\zeta|)\).
If the forcing vanishes strictly between sites \(i<j\), with gap \(d=j-i\),
and the intervening homogeneous state is nonzero, its endpoint states
are related by \(A^{d-1}\). The arithmetic-geometric mean inequality gives
\[
 \|p_{i+1}\|_1+\|p_j\|_1
                         \ge \lambda^{(d-1)/2}/32.
\]
For the internal gaps of a finite nonzero run, or all cyclic gaps of
a configuration with no zero coordinate, every endpoint coordinate
is counted at most twice. Hence
\[
 \|c\|_1\ge\frac1{128}\sum_{\text{forcing gaps }d}
                                   \lambda^{(d-1)/2}.      \tag{Ret.R5}
\]

Fix \(\alpha>0\) and define
\[
 b_\alpha=\sum_{v\ne0}e^{-\alpha\|v\|_1/2}
       =\left(\frac{1+e^{-\alpha/2}}{1-e^{-\alpha/2}}\right)^2-1,
 \quad h_\alpha(d)=e^{-\alpha\lambda^{(d-1)/2}/256},
 \quad H_\alpha(R)=\sum_{d\ge1}h_\alpha(d)R^d .
\]
The last series is finite for every finite \(R>0\).
For
\(0\le\eta\le[2b_\alpha H_\alpha(2)]^{-1}\),
\[
 \sum_{c\in\mathcal C_K}\eta^{q(c)}e^{-\alpha\|c\|_1}
 \le \exp\bigl(2b_\alpha^2H_\alpha(1)K\eta^2\bigr)
                                              +(3/4)^K.   \tag{Ret.R6}
\]

Indeed, if \(p\) has a zero vector, its nonzero coordinates split into
maximal contiguous runs. Each run's forcing support is contained in the
run and its preceding site and contains both boundary sites.
The supports for distinct runs are disjoint.
For \(q\ge2\) forcing sites and internal gaps \(d_1,\ldots,d_{q-1}\),
split the exponential weight in half. Summing one half over nonzero
forcing vectors and using (Ret.R5) for the other half gives the bound
\((\eta b_\alpha)^q\prod h_\alpha(d_i)\).
The component sum with fixed starting site is at most
\[
 \sum_{q\ge2}(\eta b_\alpha)^qH_\alpha(1)^{q-1}
                       \le2b_\alpha^2H_\alpha(1)\eta^2.
\]
Allowing arbitrary component subsets gives the exponential in (Ret.R6).

If every \(p_j\) is nonzero, mark a nonzero forcing site. For \(q\) such sites
the cyclic gaps sum to \(K\), and every unmarked configuration is counted
\(q\) times. Its total weight is bounded by
\[
 \sum_{q=1}^K\frac Kq(\eta b_\alpha)^q[z^K]H_\alpha(z)^q
 \le K2^{-K}\sum_{q\ge1}
                  \frac{(\eta b_\alpha H_\alpha(2))^q}{q}
 \le K2^{-K}\log2\le(3/4)^K .
\]
This includes every noninteger class. Counting is over forcing vectors
with unique real preimages, so no extra factor of 32 is needed.

The power \(\eta^2\) in (Ret.R6) cannot in general be replaced by \(\eta^3\):
a run with forcing \((-A^dv,v)\) at two sites exists for every \(d\ge1\).
The third power of the kick needed below will instead come from the
Fourier-atom degree after the reduction (Ret.R4).

\hypertarget{exact-replica-phases-for-vector-charges}{%
\subsection{Exact replica phases for vector charges}\label{exact-replica-phases-for-vector-charges}}

Write \(e^{i\epsilon b}=\sum_v k_v(\epsilon)e_v\).
Replacing the physical kick \(e^{i\epsilon\operatorname{Op}(b)}\) by
\(\operatorname{Op}(e^{i\epsilon b})\) costs at most
\(\pi\epsilon^2F_1(b)^2e^{|\epsilon|F_0(b)}/(2N)\) per kick.
The total \(rK\)-operator error is bounded by
\[
                \frac{C_bKs^2}{rN}e^{KsF_0(b)}.            \tag{Ret.R7}
\]
For sufficiently small fixed \(s_0\) this is \(O_b(e^{-cK})\).
The original physical operator is retained through this explicit bound;
commuting different Fourier modes within a kick would be incorrect.

Let \(v_{\nu j}\in\mathbb Z^2\) be its Fourier index in replica \(\nu\)
and at time \(j\), and \(c_j=\sum_\nu v_{\nu j}\).
Regroup the original replica-ordered word by sites, keeping replica order
within a site. Its exact crossing phase is
\[
 \exp\left[-\frac{2\pi i}{N}
     \sum_{i<j}\sum_{\nu<\mu}
                 \omega(v_{\mu i},A^{j-i}v_{\nu j})\right]. \tag{Ret.R8}
\]
The regrouped word contains the onsite factors
\[
 \exp\left[\frac{\pi i}{N}\sum_j\sum_{\nu<\mu}
                      \omega(v_{\nu j},v_{\mu j})\right]   \tag{Ret.R9}
\]
and the remaining product \(\prod_j\mathcal W_N(A^jc_j)\).
Only this remaining product depends solely on the nonzero aggregates.
Both (Ret.R8) and (Ret.R9) can depend on clean replica configurations.
The sign in (Ret.R8) follows from each inversion:
\(\mathcal W(a)\mathcal W(d) =e^{-2\pi i\omega(d,a)/N}\mathcal W(d)\mathcal W(a)\).

In (Ret.R8), replace \(j-i>m\) by \(j-i-K\). This preserves the phase because
\(A^K=I\bmod N\). Its bilinear kernel has a uniformly bounded spatial row
sum. Using the maximum matrix-entry norm, an admissible bound is
\[
 \sum_{d=1}^{K-1}\frac{2\pi}{N}
             \max|A^{\operatorname{center}(d)}|
 \le2\pi\left(4+\frac8{\lambda-1}\right)=8\pi\sqrt2.        \tag{Ret.R10}
\]
Here \(\max|A^j|=4b_{|j|}\) for \(j\ne0\).
The bound on a particle interaction is this kernel times
\(l(\boldsymbol v)l(\boldsymbol w)\), where
\(l(\boldsymbol v)=\sum_\nu\|v_\nu\|_1\), independently of \(r\).

The antipodal phase must not be set to one for general vector charges:
\(\omega((1,0),A^m(0,1))=\alpha_m\), and
\(e^{2\pi i\alpha_m/N}\ne1\). The position-only cancellation at \(d=m\)
does not apply here.

\hypertarget{the-clean-complex-gas-including-its-onsite-correction}{%
\subsection{The clean complex gas, including its onsite correction}\label{the-clean-complex-gas-including-its-onsite-correction}}

At a clean site \(\sum_\nu v_\nu=0\). The vacuum is
\(v_0=k_0(\epsilon)^r\); every nonvacuum particle has at least two nonzero
replica entries. Its activity is
\[
 w_j(\boldsymbol v)=v_0^{-1}\prod_\nu k_{v_\nu}(\epsilon)
   \exp\left[\frac{\pi i}{N}\sum_{\nu<\mu}\omega(v_\nu,v_\mu)
                         +i\sum_\nu\langle h_{\nu j},v_\nu\rangle\right].
\]
The fields are real and arise from specified dirty profiles through (Ret.R8).
The onsite factor has modulus one.
For a fixed positive analytic weight \(\beta\) smaller than the strip
width, the nonvacuum activity moment satisfies, uniformly in the fields,
\[
 \sup_j\sum_{\boldsymbol v\ne0}
           |w_j(\boldsymbol v)|\,l(\boldsymbol v)e^{\beta l(\boldsymbol v)}
                                                    \le C_bs^2. \tag{Ret.R11}
\]
This follows by summing the two required nonzero entries, using
\(\sum_{v\ne0}|k_v|e^{\beta\|v\|}=O_b(|\epsilon|)\),
the corresponding first-moment bound, and
\(|k_0(\epsilon)|^{-r}\le e^{C_br\epsilon^2}\).
The replica binomial factors are bounded by powers of \(r\), absorbed in \(s\).

Combine (Ret.R10), same-site exclusion, and (Ret.R11) with the complex tree
estimate proved in the fixed-harmonic note. Its criterion is
\((1+8\pi\sqrt2)M_\beta<\beta\).
On any set of clean sites the connected logarithm after the one-particle
term is \(O_b(Ks^4)\), uniformly in \(K,r\) and all real fields.

One must also control the single-site sum. Without (Ret.R9) it is
\[
 B^{\rm cl}_j=\int_{\mathbb T^2}
       \exp\left(i\epsilon\sum_{\nu=1}^r
                       b(z+h_{\nu j}/(2\pi))\right)\,dz .
\]
With (Ret.R9), denote it by \(B_j\). Fourier first moments give
\[
 |B_j-B^{\rm cl}_j|
       \le\frac{\pi s^2F_1(b)^2}{2N}e^{sF_0(b)}.           \tag{Ret.R12}
\]
To see this, bound the difference of the onsite phase from one by
\(\pi N^{-1}\sum_{\nu<\mu}\|v_\nu\|_1\|v_\mu\|_1\),
drop the zero-sum constraint, and use
\(\sum_v\|v\|_1|k_v|\le|\epsilon|F_1(b)e^{|\epsilon|F_0(b)}\).
The pair count is at most \(r^2/2\).

Since \(b\) has mean zero, Taylor expansion and the gas remainder imply
\[
 \log|\mathcal Q_{B,\boldsymbol h}|
 \le-\frac{\epsilon^2}{2}\sum_{j\notin B}
       \left\|\sum_\nu b(\,\cdot+h_{\nu j}/(2\pi))\right\|_2^2
                          +C_bK(s^3+s^2/N).               \tag{Ret.R13}
\]
Vacuum factors at the clean sites are included in \(\mathcal Q\).
For precision, the logarithmic gas formula is
\(\sum_{j\notin B}\log v_0+\sum_{j\notin B}(B_j/v_0-1)+O(Ks^4)\).
Because \(B_j/v_0-1=O(s^2)\), replacing its first two terms by
\(\sum\log B_j\) costs \(O(Ks^4)\). This proves (Ret.R13).
At no external field the full clean partition \(Q_{r,K}\) is nonzero and
\[
 \left|\log Q_{r,K}+\tfrac12\sigma^2K(r\epsilon)^2\right|
                                   \le C_bK(s^3+s^2/N).   \tag{Ret.R14}
\]

For dirty profiles of total absolute mass
\(S_B=\sum_{j\in B,\nu}\|v_{\nu j}\|_1\), the field row bound gives
\(\sum_{j\notin B}\max_\nu\|h_{\nu j}\|_\infty\le C_AS_B\).
Translation in \(L^2\) and (Ret.R13) consequently yield
\[
 |\mathcal Q_{B,\boldsymbol h}|
 \le \exp\left[-\tfrac12\sigma^2Ks^2+C_bK(s^3+s^2/N)
                                  +C_bs^2(|B|+S_B)\right]. \tag{Ret.R15}
\]
Indeed \(\|\sum_\nu b(\cdot+a_\nu)\|_2^2 \ge r^2\sigma^2-2r\|b\|_2\|\nabla b\|_2\sum_\nu\|a_\nu\|\).
This estimate requires no smallness of the dirty charges or their fields.

\hypertarget{why-the-finite-component-contribution-starts-at-cubic-order}{%
\subsection{Why the finite-component contribution starts at cubic order}\label{why-the-finite-component-contribution-starts-at-cubic-order}}

Let \(D>0\) dominate the constant multiplying \(s^2(|B|+S_B)\) in (Ret.R15).
Define a Fourier series with positive coefficients
\[
 B_s(z)=\sum_{v\ne0}|\widehat b_v|e^{Ds^2\|v\|_1}e_v(z),
 \qquad
 a_c(s)=e^{Ds^2}\,[e^{sB_s}]_c\quad(c\ne0).                 \tag{Ret.R16}
\]
For a dirty site, the absolute sum over all its replica profiles,
including the factor from (Ret.R15), is bounded by \(a_c(s)\).
The inequality follows by expanding each \(k_v(\epsilon)\) into its
original Fourier atoms and using the triangle inequality on their
aggregate norm. Summing replica assignments gives \(r|\epsilon|=s\).
This also accounts for internal cancellations inside a dirty replica.

For some \(\alpha>0\) and small \(s_*>0\),
\(a_c(s)\le C_bs\,e^{-\alpha\|c\|_1}\) for \(0\le s\le s_*\).
The same estimate with \(|s|\) bounds its absolute Taylor series on a
complex disk: choose \(Ds_*^2\) smaller than half the available strip width.
Each \(a_c(s)\) has nonnegative Taylor coefficients.

Let \(\mathcal M(s)\) be the sum, with one fixed starting site, of
\(\prod_{c_j\ne0}a_{c_j}(s)\) over all finite nonzero integer preimage runs.
The component proof of (Ret.R6) gives a normally convergent analytic series
\(\mathcal M(s)=O_b(s^2)\) on a fixed disk.
Its coefficient of degree two is, however, zero.
A term of that degree would use two original Fourier atoms of \(b\)
at two distinct forcing sites. Their exact lifted relation would be
\(A^iv+A^jw=0\) with \(i\ne j\).
The orbit-representative property, including the paired negative
representatives, would force \(A^{i-j}v=v\) for a nonzero integer vector.
Hyperbolicity excludes this.
Two opposite atoms at the same site have zero aggregate and do not
form a nonzero forcing component.

Since the coefficients of \(\mathcal M\) are nonnegative and its degrees
zero through two vanish, for a fixed \(0<R<s_*\),
\[
 \mathcal M(s)\le(s/R)^3\mathcal M(R)\le C_bs^3
                                      \quad(0\le s\le R). \tag{Ret.R17}
\]
This is a bound on the total original Fourier-atom degree, not a false
assertion that every vector component occupies at least three sites.
For example, the orbit-representative modes
\((1,0),(0,1),(1,1)\) allow the two-site forcing
\[
 c_0=(1,1),\qquad c_1=(1,-1),
 \qquad c_0+Ac_1=0.
\]
It uses three original atoms: one at site zero and two at site one.

The nonwinding sum over nonzero aggregates is at most
\(e^{K\mathcal M(s)}-1\). For winding configurations apply (Ret.R6) with
\(\eta=C_bs\) to get \((3/4)^K\) after reducing \(R\).
All scalarized central words and all 32 classes are therefore included.
Writing their normalized full trace as \(\mathcal T_{r,K}\), (Ret.R15) gives
\[
 |\mathcal T_{r,K}-Q_{r,K}|
 \le e^{-\sigma^2Ks^2/2+C_bK(s^3+s^2/N)}
                 \left[e^{C_bKs^3}-1+(3/4)^K\right].       \tag{Ret.R18}
\]
No factorization of the quantum phases of distinct dirty components
is assumed; the sum is bounded after taking absolute values.

For sufficiently large \(K\) and sufficiently small fixed \(s_0\),
the cubic and \(N^{-1}\) terms are absorbed into a fixed fraction of
the negative quadratic term. Equations (Ret.R14) and (Ret.R18) imply
\[
 |\mathcal T_{r,K}|\le e^{-\sigma^2Ks^2/4}+(3/4)^K,
\]
\[
 |\mathcal T_{r,K}-e^{-\sigma^2Ks^2/2}|
       \le C_b(Ks^3+Ks^2/N)e^{-c_bKs^2}+(3/4)^K.           \tag{Ret.R19}
\]
The maximum of \(Ks^3e^{-c_bKs^2}\) is \(O_b(K^{-1/2})\); the other
term is \(O_b(N^{-1})\). Restore the physical kick by (Ret.R7) and \(f\)
by (Ret.R4), reducing \(s_0\) so that \(s_0F_0(b)<\log\lambda/4\).
The odd correction, when relevant, is at most
\(e^{KsF_0(b)}/(N-1)\) before the already controlled replacements.
All these errors are exponentially small in \(K\).
Enlarging \(C\) handles the finitely many remaining small \(K\).
This proves (Ret.R1)--(Ret.R2).

The spectral consequences in Section 1 follow by testing each fixed
integer Fourier moment. Necessity of \(K\epsilon_K^2\to\infty\) for
Haar measure follows already from the first moment and (Ret.R2).
This proof of the analytic return law does not import a classical CLT.

\hypertarget{a-logarithmic-fourier-bound-for-quantum-period-averages}{%
\subsection{A logarithmic Fourier bound for quantum period averages}\label{a-logarithmic-fourier-bound-for-quantum-period-averages}}

For an absolutely summable Fourier series of mean zero define
\[
 X_K(f)=\frac1{\sqrt K}\sum_{j=0}^{K-1}
                U_0^j\operatorname{Op}_N(f)U_0^{-j}.
\]
This section concerns these averaged observables. It does not lower the
regularity assumption for the actual kicked propagator in (Ret.R1).

If a nonzero integer vector \(v\) has an early modular return
\(A^dv=v\bmod N\), \(0<d<K\), consider the forcing \(c_0=-v,c_d=v\).
Its preimage has no zero coordinate. Otherwise it would be integer
and would split into finite components. The only two forcing atoms
would have to belong to the same component, and its lifted exact
relation would imply \(A^ev=v\) for some nonzero \(|e|<K\).
This is impossible for a hyperbolic \(A\).
Apply (Ret.R5) to the two cyclic gaps \(d,K-d\) to obtain
\[
 2\|v\|_1\ge\frac1{128}
       \left(\lambda^{(d-1)/2}+\lambda^{(K-d-1)/2}\right),
 \qquad
 \|v\|_1\ge\frac{\lambda^{K/4}}{256\sqrt\lambda}.           \tag{Ret.R20}
\]
Thus an early return is possible only if
\(K\le C_A\log(2+\|v\|_1)\).

For one Fourier character, if all \(K\) orbit residues are distinct,
Weyl orthogonality gives \(\|X_K(e_v)\|_{2,N}=1\) exactly.
If any two coincide, use (Ret.R20) and the elementary bound
\(\|X_K(e_v)\|_{2,N}\le\sqrt K\).
Consequently, for every \(K\) and every \(v\ne0\),
\[
 \|X_K(e_v)\|_{2,N}
                      \le C_A\sqrt{\log(2+\|v\|_1)}.       \tag{Ret.R21}
\]
This includes central aliases and possible balanced signs among repeated
residues; the upper bound does not treat all such signs as positive.
By the triangle inequality,
\[
 \sup_K\|X_K(f)\|_{2,N}\le C_A\mathcal L_{1/2}(f),\qquad
 \mathcal L_\beta(f)=
       \sum_{v\ne0}|\widehat f_v|[\log(2+\|v\|_1)]^\beta.  \tag{Ret.R22}
\]

\textbf{Theorem 2 (low-regularity averaged quantum CLT).} If \(f\) is real,
mean zero and \(\mathcal L_{1/2}(f)<\infty\), the empirical spectral law
of \(X_K(f)\) converges weakly to \(N(0,\sigma_A^2(f))\).
More generally, for every fixed finite family in this class and fixed
real \(t_1,\ldots,t_J\),
\[
 \frac1N\operatorname{Tr}\prod_{j=1}^J e^{it_jX_K(f_j)}
       \longrightarrow
       \exp\left[-\tfrac12\sigma_A^2\left(\sum_jt_jf_j\right)\right].
                                                               \tag{Ret.R23}
\]
For even symbols the statement holds on the odd block as well.
No convergence of arbitrary unbounded moments is asserted.

For a fixed Fourier polynomial, (Ret.R2), or (Ret.R3) when its variance vanishes,
and the gauge-free estimate (Ret.P31) give the scalar statement.
For finitely many such polynomials,
(Ret.P32) and iterated unitary Duhamel give the joint statement.
These uses of (Ret.P31)--(Ret.P32) require only their finite-dimensional estimates,
not the classical CLT imported elsewhere in the companion note.

Now let \(f_M\) be a fixed finite Fourier cutoff. By (Ret.R22),
\[
 \sup_K\|X_K(f)-X_K(f_M)\|_{2,N}
                          \le C_A\mathcal L_{1/2}(f-f_M)\to0.
\]
The corresponding exponential difference is bounded by \(|t|\) times
this norm. Telescope a fixed product of unitary exponentials to obtain
the same approximation for its normalized trace.
First take \(K\to\infty\) at fixed cutoffs, then let the cutoffs grow.
The covariance converges by absolute Fourier summability and the
orbit formula (Ret.P18). This proves (Ret.R23) and, by characteristic functions,
the weak spectral convergence. The transfer to the odd space costs
the bounded norm factor \(\sqrt{2N/(N-1)}\).

\hypertarget{the-exponent-one-half-is-sharp-within-these-weighted-fourier-classes}{%
\subsection{The exponent one half is sharp within these weighted Fourier classes}\label{the-exponent-one-half-is-sharp-within-these-weighted-fourier-classes}}

\textbf{Theorem 3.} There is a single fixed real even continuous mean-zero
function \(f\) for which \(\mathcal L_\beta(f)<\infty\) for every
\(\beta<1/2\), but the empirical laws of \(X_K(f)\) are not tight.
In particular replacing \(1/2\) by any smaller exponent in Theorem 2
would be false. This is sharpness of a uniform class guarantee, not
a necessary condition for every individual observable.

For \(j\ge1\) put
\[
 K_j=4\cdot2^{\,2^j}+2,\qquad
 N_j=b_{K_j/2},\qquad c_j=\frac j{\sqrt{K_j}},\qquad
 f(z)=2\sum_{j\ge1}c_j\cos(2\pi N_jz_1).                  \tag{Ret.R24}
\]
Every \(K_j\) is admissible. The coefficients are summable, so this
defines one fixed continuous function, not a dimension-dependent symbol.
Since \(\log N_j\asymp K_j\),
\[
 \mathcal L_\beta(f)
      \le C_\beta\sum_j j K_j^{\beta-1/2}<\infty
                                          \quad(\beta<1/2).
\]
At dimension \(N_j\), the \(j\)th character is the identity:
\(\mathcal W_{N_j}(N_j,0)=I\). It stays the identity under every
Egorov iterate. Its contribution to \(X_{K_j}(f)\) is therefore \(2jI\).

Every earlier frequency \(N_i(1,0)\), \(i<j\), has \(K_j\) distinct orbit
residues for all sufficiently large \(j\). Indeed \(N_i<\lambda^{K_i/2}\),
whereas an early return would require the lower bound (Ret.R20).
The difference \(K_j-2K_{j-1}\) tends to infinity.
Thus the contribution \(Y_j\) of the earlier modes has
\[
                         \|Y_j\|_{2,N_j}\le2\sum_i c_i<\infty .
\]
The later modes have operator norm at most
\(2\sqrt{K_j}\sum_{i>j}c_i\).
The recurrence \(K_{j+1}=(K_j-2)^2/4+2\) gives
\(K_{j+1}\ge K_j^2/5\); for \(j\ge2\) it also gives
\(c_{j+1}/c_j\le1/2\).
Hence their contribution \(E_j\) satisfies
\[
 \|E_j\|_{\rm op}
       \le4(j+1)\sqrt{K_j/K_{j+1}}
       \le\frac{4\sqrt5(j+1)}{\sqrt{K_j}}\longrightarrow0,
 \qquad X_{K_j}(f)=2jI+Y_j+E_j.                            \tag{Ret.R25}
\]
The spectral Chebyshev inequality and the Hermitian min-max principle
now show that the fraction of eigenvalues of \(X_{K_j}(f)\) below \(j\)
is at most \(C/(j-\|E_j\|_{\rm op})^2\), which tends to zero.
The empirical spectra escape to positive infinity along this subsequence.
The same argument works on the odd block.

The classical variance is nevertheless finite and nonzero:
\(\sigma_A^2(f)=2\sum_jc_j^2\).
Different positive multiples of \((1,0)\) belong to different oriented
\(A\)-orbits, and nonzero-time classical correlations vanish.
The failure comes from quantum aliases of the high Fourier modes.
This continuous example is outside the smooth and analytic kick classes;
it is not a counterexample to their return theorems or to a formulated
fixed-system BGS conjecture.

\hypertarget{the-failure-persists-after-subtracting-the-quantum-mean}{%
\subsection{The failure persists after subtracting the quantum mean}\label{the-failure-persists-after-subtracting-the-quantum-mean}}

The previous example uses a divergent scalar alias. The loss of tightness
below the exponent \(1/2\) is not repaired by subtracting that scalar:
there is also one fixed real even continuous function \(f\) in every
class \(\mathcal L_\beta\), \(\beta<1/2\), for which
\[
 \widetilde X_K(f)=X_K(f)-\frac1N\operatorname{Tr}X_K(f)\,I
                                                               \tag{Ret.R26}
\]
has non-tight spectral laws.

Here is an explicit construction. Use the \(K_j\) of (Ret.R24) and set
\[
 L_j=15K_j,\quad M_j=b_{L_j/2},\quad d_j=M_j/5,\quad
 a_j=j/\sqrt{L_j},\qquad
 f(z)=2\sum_{j\ge1}a_j\cos(2\pi d_jz_1).                    \tag{Ret.R27}
\]
The same summability argument proves the stated regularity.
All \(L_j\) are admissible. The exact recurrence gives \(25\mid b_{15}\).
Since \(15\mid L_j/2\) and \(b_u\mid b_{uv}\) (apply the formula for
\(A^{uv}\) to \(A^u=\alpha_uI+b_uB\)), we have \(25\mid M_j\).

Fix one such dimension \(M=25t\), where \(t\) is odd, and let \(d=M/5\).
On the physical \(M\)-dimensional space write
\[
 C=\mathcal W_M(d,0),\qquad S=\mathcal W_M(0,d).
\]
They commute, have order five, and
\[
 \mathcal W_M(dv)=(-1)^{v_1v_2}C^{v_1}S^{v_2}.
\]
Indeed their commutator is \(e^{2\pi i M/25}=1\), and the balanced
onsite factor is \(e^{-\pi i(M/25)v_1v_2}=(-1)^{v_1v_2}\).
Their 25 joint characters have equal multiplicity \(M/25\):
the normalized trace of \(C^uS^v\) is zero unless \(u=v=0\bmod5\).

Since \(A=I\bmod2\), every \(A^\ell(1,0)\) has even coordinate product.
Modulo five its orbit has length six and consists of
\[
 (1,0),\ (3,2),\ (2,2),\ (-1,0),\ (-3,-2),\ (-2,-2).
\]
Because \(6\mid L_j\), the averaged character at frequency \(d_j(1,0)\) is
the Hermitian contraction
\[
 H_j=\frac16\sum_{\ell=0}^5
       C^{(A^\ell(1,0))_1}S^{(A^\ell(1,0))_2},\qquad
 \frac1{M_j}\operatorname{Tr}H_j=0.                        \tag{Ret.R28}
\]
It equals the identity on the joint \(C=S=1\) subspace of dimension
\(M_j/25\). Thus the contribution of the \(j\)th real Fourier mode to
\(X_{L_j}(f)\) is \(2jH_j\), rather than a scalar matrix.

The earlier frequencies have full modular period \(L_j\) for all
sufficiently large \(j\), by (Ret.R20). Their contribution \(Y_j\) has
uniformly bounded normalized Hilbert--Schmidt norm and zero trace.
The later frequencies contribute an operator \(E_j\) with norm tending
to zero, by the same geometric-tail estimate as (Ret.R25). Hence
\[
 \widetilde X_{L_j}(f)=2jH_j+D_j,\qquad
 D_j=Y_j+E_j-\frac1{M_j}\operatorname{Tr}E_j\,I,\qquad
 \sup_j\|D_j\|_{2,M_j}<\infty.                             \tag{Ret.R29}
\]
The spectral subspace where \(|D_j|>j\) has dimension at most
\(M_j\|D_j\|_{2,M_j}^2/j^2\). Its orthogonal complement intersects the
\(H_j=1\) subspace in dimension at least the difference of these ranks.
On this intersection the quadratic form of \(\widetilde X_{L_j}(f)\)
is at least \(j\). The min-max principle therefore gives
\[
 \frac1{M_j}\#\{\text{eigenvalues of }\widetilde X_{L_j}(f)
                                           \text{ at least }j\}
       \ge\frac1{25}-\frac{C_f}{j^2}.                     \tag{Ret.R30}
\]
This proves non-tightness even with the actual quantum mean removed.

The conclusion also holds after restriction to the odd space and
subtraction of its own normalized trace. To count the relevant
subspace, an orthonormal basis of \(C=S=1\) is
\[
 w_a=\frac1{\sqrt5}\sum_{\ell=0}^4|5a+\ell d\rangle,
 \qquad a\in\mathbb Z/t\mathbb Z .
\]
Parity sends \(w_a\) to \(w_{-a}\), so its odd dimension is \((t-1)/2\).
Its relative dimension in the odd space is
\((M/25-1)/(M-1)\to1/25\).
The normalized odd trace of \(H_j\) is \(O(M_j^{-1})\), since the
parity-twisted trace of each Weyl operator has modulus one.
Subtracting this trace changes \(2jH_j\) by \(O(j/M_j)=o(1)\).
The odd Hilbert--Schmidt norm of \(Y_j\) is bounded; subtracting its
odd trace preserves a uniform bound.
The same intersection and min-max argument, with threshold \(j/2\) if
necessary, gives a positive limiting fraction of escaping eigenvalues.

\hypertarget{verification-and-scope}{%
\subsection{Verification and scope}\label{verification-and-scope}}

The companion verifier is
\href{https://github.com/ChurilovMaksim/BGS/blob/main/computations/verification/phase_space_return_gas.py}{phase\_space\_return\_gas.py}.
Run from the repository root with NumPy available:

\begin{verbatim}
OPENBLAS_NUM_THREADS=1 python computations/verification/phase_space_return_gas.py \
  --output build/phase-space-return-gas-precheck.json
\end{verbatim}

The recorded run passes the following controls:

\begin{itemize}
\tightlist
\item
  320 exact replica-word reordering identities at \(K=6,10,14,34\);
  explicit nontrivial antipodal and onsite phases;
\item
  3,000 rational endpoint inequalities, checked in
  \(\mathbb Q(\sqrt2)\) without floating point, and 120 finite
  two-forcing-site runs;
\item
  1,224 finite two-atom exclusions and an explicit three-atom component
  occupying only two forcing sites;
\item
  25 early-return controls through \(K=262\), with the stronger two-gap
  inequality checked exactly when a return occurs;
\item
  exact recurrence and tail-ratio checks, central-alias controls at
  \(K=18,66,1026\), and noncentral period-six aliases at
  \(L=270,990,15390\). Older modes have the full target periods in
  the checked cases; the proof above establishes the eventual assertion;
\item
  six finite Taylor-polynomial onsite matrix comparisons at \(N=35\),
  separately labeled as floating diagnostics. Commuting subgroup
  projectors at \(M=25,75,125\) reproduce their full and odd ranks;
  the largest matrix residual is below \(4.4\cdot10^{-14}\).
\end{itemize}

The analytic lattice, conditional gas, cubic-degree mechanism, and both
regularity counterexamples received limited independent internal review.
That review led to two clarifications: the forcing support is
contained in a run and its preceding site, and the zero-variance
case of Theorem 2 uses (Ret.R3). This is not external referee acceptance
or formal verification.

The actual-kick theorem requires analytic \(f\), while the low-regularity
theorem concerns \(X_K(f)\). Uniformity in \(r\) in Theorem 1 requires
\(r|\epsilon|\le s_0\): it does not permit arbitrary time growth at
a fixed kick. Neither global Gaussian return laws nor a regularity
counterexample supply local spectral correlations for a fixed
nonzero classical perturbation.

The subsequent \href{https://github.com/ChurilovMaksim/BGS/blob/main/mathematics/general_cat_return_laws.md}{general cat-map argument}
extends the analytic return theorem, the logarithmic regularity
threshold and its centered counterexample to every fixed hyperbolic
\(A\in SL(2,\mathbb Z)\), on explicitly specified arithmetic subsequences.
The Gaussian CLT for classical toral sums is prior mathematics.
Its source-level check is in
\href{https://github.com/ChurilovMaksim/BGS/blob/main/literature/periodic_return_clt_sources.md}{the literature note}.
No claim of a first result in the literature follows from these
project derivations.
